\section{Extended Proofs}
\label{app:proofs}

This appendix collects proofs that are deferred from the main text
for readability, and provides extended versions of proofs given in
sketch form.

\subsection{Proofs of Propositions~\ref{prop:channel_isolation}
and~\ref{prop:channel_independence}
(Channel Isolation and Independence)}
\label{app:channel_independence}

The main-text proofs (\S\ref{sec:channel_independence}) are now the
primary versions.  Proposition~\ref{prop:channel_isolation} (Channel
Isolation) provides the result that is robust against single-substrate
coalitions, requiring only substrate exclusivity;
Proposition~\ref{prop:channel_independence}
(Channel Independence) strengthens to full conditional independence
under additional joint-sufficiency and independent-randomness
assumptions.
We provide supplementary discussion of the conditioning requirements
for the stronger result.

\paragraph{Why condition on benchmark parameters.}
The governance process that accepts benchmarks $B_1, \ldots, B_n$ is
a shared cause: if a single governance round accepts correlated
benchmarks (e.g., all benchmarks drawn from the same test suite family),
verdicts can be correlated through shared difficulty even without
witness cross-communication.  Conditioning on the benchmark specifications
removes this path.  The proposition's requirement that benchmarks be
``independently specified by governance'' is an assumption, not a
consequence of the governance fork protocol.  Separate governance
proposals and votes are a necessary procedural condition (each
benchmark goes through its own deliberation), but they do not
guarantee statistical independence of benchmark parameters; for
example, benchmarks proposed by the same entity or drawn from the
same test-suite family may share structural biases.  The assumption
holds when the governance process produces genuinely independent
benchmark designs; the proposition's guarantee is conditional on it.

\paragraph{Why condition on capability profile (joint sufficiency).}
Agent $a$'s internal state (training history, memory, strategic
disposition) is a shared cause across all exercises.  The joint
sufficiency assumption (iii) requires that $\mathbf{C}_a$ absorbs
all such cross-exercise dependencies: conditional on $\mathbf{C}_a$,
the agent's performance on $\pi_{B_i}$ is independent of its
performance on $\pi_{B_j}$.  This is strictly stronger than marginal
sufficiency (which would only require that $\mathbf{C}_a$ predicts
per-benchmark performance without guaranteeing cross-benchmark
independence).  The joint requirement excludes persistent
cross-exercise state (e.g., an agent that learns from one exercise
and adapts its strategy for the next) and shared agent-side
randomness not absorbed by $\mathbf{C}_a$.
If the agent has latent strategic state not captured by $\mathbf{C}_a$
(e.g., a deceptive agent that adapts its behavior per-benchmark),
the conditional-independence guarantee of
Proposition~\ref{prop:channel_independence} weakens, but the
channel-isolation guarantee of
Proposition~\ref{prop:channel_isolation} still holds, because the
isolation result depends only on substrate exclusivity, not on the
agent's internal state.  This is the structural reason for separating
the two results: isolation is the load-bearing property against
single-substrate coalitions, and independence is the additional
statistical guarantee available when the agent's capability profile
is a jointly sufficient summary of its exercise-level behavior.

\subsection{Proof of Proposition~\ref{prop:coalition_resistance}
(Coalition Resistance)}
\label{app:coalition_resistance}

\begin{proof}
Let $K = \{a_1, \ldots, a_k\}$ be a coalition of agents on substrate
$s$, and let $\Witness_{\mathrm{alg}}$ be an algorithmic witness on
substrate $s' \neq s$ executing protocol $\pi_B$.

The verification output is $v = \pi_B(a, \Witness_{\mathrm{alg}})$,
where $\pi_B$ is fully specified by the governance-accepted benchmark
$B$ (Definition~\ref{def:alg_witness}, condition~2).

\paragraph{Attack surface analysis.}
The coalition $K$ can attempt to influence $v$ through three channels:
\begin{enumerate}
    \item \textbf{Modifying the witness's computation.}  This requires
    writing to $\Witness_{\mathrm{alg}}$'s inputs on substrate $s'$.
    Since $K$ is confined to substrate $s$ and $s \neq s'$, this
    requires crossing the substrate boundary, i.e., compromising $s'$
    (case~a of the proposition).
    \item \textbf{Modifying the protocol.}  Protocol $\pi_B$ is
    specified by governance-accepted benchmark $B$.  Changing $B$
    requires the governance fork protocol
    (Definition~\ref{def:gov_fork}), which demands cross-substrate
    supermajority.  A coalition confined to substrate $s$ cannot achieve
    cross-substrate supermajority without validators on $s' \neq s$
    (case~b).
    \item \textbf{Influencing the agent's behavior during the exercise.}
    The isolation condition of Definition~\ref{def:exercise_protocol}
    requires that the agent interact only with the witness through the
    protocol-specified interface during the exercise.  Coalition
    assistance (pre-computed answers, same-substrate coaching) violates
    this condition and is excluded by assumption.  If isolation is
    imperfect---the agent receives some external signal despite the
    protocol's isolation requirement---the evaluated object is no
    longer the individual agent but a coalition-assisted system, and
    the soundness parameter $\epsilon$ (calibrated for individual
    agents) does not directly bound the assisted system's false-pass
    probability.  The defense against imperfect isolation is therefore
    not soundness per se, but \emph{detection and re-testing}: if
    coalition assistance is suspected, governance can mandate
    re-verification under tighter isolation controls (e.g., network
    sandboxing, monitored execution environments).  The proposition's
    guarantee is conditional on isolation holding; imperfect isolation
    degrades the guarantee proportionally to the assistance received.
\end{enumerate}

Therefore, the coalition's only viable attack vectors require breaching
the cross-substrate boundary (case~a) or the governance threshold
(case~b).  Within-substrate actions cannot influence the witness's
verdict beyond the soundness bound $\epsilon$.
\end{proof}

\subsection{Proof of Theorem~\ref{thm:anti_monopolar}
(Verification Impossibility under Monopolarity)}
\label{app:anti_monopolar}

The proof in \S\ref{sec:anti_monopolar} is already complete.  We
provide an extended discussion of the theorem's consequences, the
role of its physical premise, and the computational action-model
assumption.

\paragraph{The physical premise.}
Same-substrate addressability (Definition~\ref{def:same_substrate})
is the load-bearing assumption.  Without it, the theorem would be a
definition chase: ``we required $s' \neq s$; if $m = 1$, no such $s'$
exists.'' With it, the theorem says something about the physical
world: processes sharing a computational medium have nonzero causal
access to each other, and this access cannot be eliminated by any
computation \emph{within} that medium.  The impossibility is
grounded in the physics of shared substrates, not in the definitions
of this paper.

\paragraph{Relationship to the GFM anti-monopolar result.}
Paper~3~\citep{lasser2026horizon} proves that monopolar convergence
($m \to 1$) is anti-maximizing: $\volL$ decreases because
cross-substrate cooperative capabilities are destroyed.
Theorem~\ref{thm:anti_monopolar} establishes a complementary result:
monopolar convergence makes the \emph{loss of $\volL$}
\emph{unverifiable from within the population}.

The two results compose multiplicatively.  \citet{lasser2026horizon} says: ``if you
become monopolar, you lose the safety guarantee.''
Theorem~\ref{thm:anti_monopolar} says: ``if you become monopolar, you
also lose the ability to detect that you have lost it.''  The compound
failure is strictly worse than either result alone, because it closes
the self-correction loop: a population that could detect its loss of
safety could in principle take corrective action (increasing substrate
diversity), but a population that cannot detect the loss has no signal
to act on.

\paragraph{The computational action-model assumption.}
Part~(b) relies on the assumption that the population's available
actions are computations on currently occupied substrates.  This
excludes physical-world interventions such as procuring new hardware
or recruiting biological observers.  The exclusion is deliberate:
the theorem characterizes what is achievable \emph{computationally},
not what is achievable with arbitrary physical-world agency.  A
population that provisions a new substrate has changed $m$ from~1
to~2 and is no longer monopolar, so the theorem's premise no longer
holds.  The action-model restriction is therefore not a weakness but
a precise scoping of the impossibility: within the computational
resources of a monopolar population, exogenous verification cannot
be constructed.

\paragraph{Tightness.}
The bound is tight in the sense that $m = 2$ is sufficient for at
least one substrate-exclusive witness to exist.  With $m = 2$, agents
on substrate $S_1$ can be witnessed by processes on $S_2$ and vice
versa.  The number of independent verification channels scales as
$m - 1$ (Remark after Proposition~\ref{prop:channel_independence}),
so verification depth increases with substrate diversity, but the
qualitative threshold (existence versus non-existence of exogenous
verification) is at $m = 2$.

\paragraph{Corollary~\ref{cor:monopolar_undetect} as an epistemic
constraint.}
The corollary is an epistemic impossibility result, not merely a
practical difficulty.  It states that within a monopolar population,
the \emph{concept} of exogenous verification is empty: there is no
process available to the population that could serve as an independent
check.  The population may believe its verification is sound (all
evaluators report success), but this belief is unfalsifiable from
within---precisely because every potential falsifier shares the
substrate that the evaluated agents can influence.
